\section{Calibration, implementation error, and the evidence boundary}
\label{sec:deploymentcalibration}

This section supplies a finite, explicit interface between a fitted event model
and an implemented policy. The inequalities are elementary finite-horizon
perturbation and concentration results; their purpose is to make the
assumptions and units of a deployment claim auditable, not to claim new
concentration theory.

\subsection{An observed decision process and its clocks}

Let $k=0,\ldots,N$ index decision opportunities. A state $x$ records the
available public history or a declared sufficient summary, private order
reports, outstanding commitments, clock time and remaining horizon.
A report belongs to the controller's information only after receipt and
sequence recovery. An exchange timestamp is not a receipt timestamp.
For genuinely hidden state, $x$ must be a posterior or an observation history;
replacing it by a point estimate needs an additional reduction allowance.
An event-time model must retain holding times or calendar time whenever
costs, signal decay or the terminal deadline depend on them. If the number
of events before a deadline is unbounded, the finite $N$ result below applies
to a declared truncation or batching scheme, with an additional tail error.

A feasible action $a\in A_k(x)$ is a message or a deliberately empty message.
Write $P_k^a(x,\cdot)$ for the next-state kernel and $r_k(x,a)$ for the
conditional mean net reward accrued before the next decision. Partial fills,
fees, simultaneous price changes and elapsed-time charges are jointly
accounted before taking this mean. For execution, reward is negative cost.
The additive objective is
\[
 J(\pi)=\E^\pi\!\left[\sum_{k=0}^{N-1}r_k(X_k,A_k)+g(X_N)\right].
\]
This definition does not identify risk-penalized reward with cash profit.
A realized random reward can replace $r_k$ inside the expectation.

All comparisons start from the same fixed initial state or initial law.
The model and deployment use the same observation class and feasible sets.
A support error that admits overfills or a missing action that changes
feasibility is not repaired by a small reward-error bound. A conservative
common feasible mask is permitted; the resulting benchmark is the optimum
within that mask. Comparison with a larger action class needs a separate
restriction gap.

\subsection{A finite-horizon transfer with a compiled policy}

Let hats denote the fitted model. Total variation uses the convention
$\operatorname{TV}(p,q)=\tfrac12\sum_y|p_y-q_y|$.
Suppose, uniformly over feasible states and actions,
\[
 |r_k-\widehat r_k|\le\eta_k,\quad
 \operatorname{TV}(P_k^a,\widehat P_k^a)\le\epsilon_k,
 \quad |g-\widehat g|\le\eta_N.
\]
Let $S_k$ bound the span of every fitted-model continuation value of a
policy in the comparison class, including its optimal value. Bounded
stage rewards give the simple choice
$S_k\le\sum_{j=k}^{N-1}R_j+R_N$, where $R_j$ is a common range width
across all states and actions of fitted stage reward and $R_N$ is the
terminal range width. Set
\begin{equation}\label{eq:deployB}
 B_N=\eta_N,\qquad
 B_k=B_{k+1}+\eta_k+\epsilon_k S_{k+1}.
\end{equation}

\begin{theorem}[Statistical and implementation errors in one objective]
\label{thm:deploymenttransfer}
Under these assumptions, $|J(\pi)-\widehat J(\pi)|\le B_0$ for each
common feasible policy, and $|J^*-\widehat J^*|\le B_0$.
Suppose the actual implemented policy $\widetilde\pi$ satisfies
\[
 \widehat V_k^*(x)-\widehat Q_k^*(x,\widetilde\pi_k(x))
 \le\zeta_k(x),\qquad \zeta_k(x)\ge0,
\]
where $\widehat Q_k^*=\widehat r_k+\widehat P_k\widehat V_{k+1}^*$.
Then
\begin{equation}\label{eq:deploymentregret}
 J^*-J(\widetilde\pi)
 \le 2B_0+
 \widehat\E^{\widetilde\pi}\sum_{k=0}^{N-1}\zeta_k(X_k).
\end{equation}
The expectation on the right is under the fitted deployed-policy law.
It may be bounded by $\sum_k\sup_x\zeta_k(x)$.
\end{theorem}
\begin{proof}
For a fixed action and continuation $v$,
$|(P-\widehat P)v|\le\operatorname{TV}(P,\widehat P)
(\max v-\min v)$. Insert and subtract the fitted continuation in the
one-step recursion. Backward induction gives \eqref{eq:deployB} for a
fixed policy and, since the action sets agree, for the optimal Bellman
operators. Telescoping the fitted optimal value along the implemented
policy gives
$\widehat J^*-\widehat J(\widetilde\pi)
=\widehat\E^{\widetilde\pi}\sum_k
[\widehat V_k^*(X_k)-\widehat Q_k^*(X_k,\widetilde\pi_k(X_k))]$.
Insert the two model-comparison bounds to obtain
\eqref{eq:deploymentregret}. The same proof applies to history states.
\end{proof}

\begin{remarkplain}[Relation to the path-coupling certificate]
When $x$ is the full structural state of Section~29 and the decision-time maps
agree, Theorem~\ref{thm:bookcoupling} applied to one decision interval
supplies the transition budget: if the one-clock mark laws differ by at most
$\eta$ in total variation, then
$\epsilon_k\le 1-e^{-\Lambda\eta(t_{k+1}-t_k)}$ in \eqref{eq:deployB}. Neither certificate
dominates the other. Corollary~\ref{cor:bookcertificate} charges one loss range $L$
against the whole-path mismatch $1-e^{-\Lambda\eta T}$, whereas \eqref{eq:deployB} weights
each interval's mismatch by the span $S_{k+1}$ of its own remaining
continuation and additionally accepts reward-mean errors $\eta_k$ and
history states.
\end{remarkplain}

If a fixed pair of baseline and candidate policies has fitted improvement
$\widehat\Delta$, its true improvement is at least
$\widehat\Delta-B_0^{\rm candidate}-B_0^{\rm baseline}$; $2B_0$ is a
common-envelope choice. This certificate may be conservative. A shared-law
advantage comparison can be tighter because common errors cancel.
The implementation allowance must evaluate the action actually emitted
after rounding, risk checks and timeouts. An optimizer tolerance before
those transformations is not an implementation guarantee.

\subsection{A transparent finite-sample construction}

\begin{proposition}[Simultaneous coverage of a finite calibration table]
\label{prop:deploymentcoverage}
Consider $L$ prespecified state--action--time rows, each with at most $M$
next-state categories. In row $\ell$, collect a fixed number $n_\ell\ge1$
of independent draws from that row's joint reward/next-state law.
The reward has a known range width $R_\ell$. Independence between reward
and next state, or between different rows, is not needed. Let empirical
means and frequencies define $\widehat r_\ell,\widehat P_\ell$.
For $0<\delta<1$, define
\begin{align}
 e_\ell&=\min\left\{1,
 \sqrt{\frac{(M+2)\log2+\log(L/\delta)}{2n_\ell}}\right\},
 \label{eq:deployeps}\\
 d_\ell&=R_\ell\min\left\{1,
 \sqrt{\frac{\log(4L/\delta)}{2n_\ell}}\right\}.
 \label{eq:deployeta}
\end{align}
With probability at least $1-\delta$, all rows satisfy
$\operatorname{TV}(P_\ell,\widehat P_\ell)\le e_\ell$ and
$|r_\ell-\widehat r_\ell|\le d_\ell$ simultaneously.
Consequently Theorem~\ref{thm:deploymenttransfer} holds on that same event,
even for a policy selected using the table.
\end{proposition}
\begin{proof}
For a fixed subset $C$ of next-state categories, the empirical indicator
mean differs from $P_\ell(C)$ by more than $u$ with probability at most
$2e^{-2n_\ell u^2}$. There are at most $2^M$ subsets; total variation is
the maximum of these absolute differences. A union bound over rows gives
$L2^{M+1}e^{-2n_\ell u_\ell^2}\le\delta/2$ with the unclipped rowwise
choice. The bounded-variable inequality gives a further failure
probability at most $\delta/2$ for the reward means. The deterministic
bounds $1$ and $R_\ell$ justify clipping. On their simultaneous event the
preceding theorem applies uniformly, so subsequent policy selection does
not require another union bound over policies.
\end{proof}

This deliberately conservative construction is a complete positive result
for a stated sampling experiment. It is not an assertion that adjacent
market events are independent. A single dependent book path, data-dependent
bins, adaptive sample counts or hidden confounding do not satisfy it merely
because many messages were recorded. Such settings require their own
martingale, mixing or block-sampling argument. A row with no supported
observations receives the trivial bounds $e_\ell=1$ and $d_\ell=R_\ell$,
or its action is excluded; it does not receive zero error.

If the deployment row is within total variation $v_\ell$ of its calibration
law and its mean reward has changed by at most $w_\ell$, use
$\epsilon_\ell=\min(1,e_\ell+v_\ell)$ and
$\eta_\ell=d_\ell+w_\ell$. This follows by the triangle inequality.
The drift budgets must be supplied or validated; a regime detector alone
does not prove them. A latency-tail, unmodeled-event or state-reduction
allowance can enter the same budgets only after its corresponding coupling
or mean-error argument is stated. There is no blanket independent-delay
assumption in this finite-kernel result.

\subsection{Replay, interventions and supported policy evaluation}

Public order-by-order data support displayed-book reconstruction. They do
not reveal all latent liquidity: the Nasdaq ITCH specification distinguishes
non-displayed executions from displayed-order lifecycle events
\citep{itch}. A simulation of an inserted order additionally chooses
its arrival time, priority and the market's response. Replaying unchanged
external messages is a small-intervention modeling assumption, not an
observed response to the inserted order. Raw packet replay helps reproduce
transport and decoder behavior; a correctly ordered event stream can be
sufficient for an economic experiment that explicitly supplies transport
delays. Neither format creates missing counterfactual data.

A simple positive off-policy statement makes the missing assumptions
precise. Consider independent one-decision records $(X_i,A_i,Y_i)$ from a
logging policy $\beta$, a candidate policy $\pi$ fixed independently of
these evaluation records, and $|Y_i|\le M_Y$. Suppose the conditional
outcome law given $(X,A)$ is invariant under the intervention, actions
are conditionally randomized or unconfounded given $X$, and
$\beta(a\mid x)\ge b_0>0$ whenever $\pi(a\mid x)>0$.
The importance-weighted estimate
\[
 \widehat J_{\rm IS}=\frac1n\sum_i
 \frac{\pi(A_i\mid X_i)}{\beta(A_i\mid X_i)}Y_i
\]
is unbiased for the candidate value on the logging context distribution.
Conditioning on $X$, summing over $A$ cancels $\beta$ and proves this claim.
Each summand lies in $[-M_Y/b_0,M_Y/b_0]$, so the bounded-variable inequality
gives the fixed-sample confidence radius
\[
 \frac{M_Y}{b_0}\sqrt{\frac{2\log(2/\delta)}{n}}.
\]
This elementary instance illustrates the established high-confidence
off-policy evaluation programme \citep{DeployThomas2015}. For sequential
policies one generally needs trajectory or sequential density ratios and
support for the induced histories; one-step weights are insufficient.
An omitted confounder or an opponent response that changes with deployment
invalidates the invariance assumption. At an unlogged action, two laws
can agree on every observed record and assign opposite candidate rewards.
No amount of additional data from the same unsupported policy identifies
that comparison.
